\problemname{Spinning Poles}
\textit{NOTE!!! This problem is not about the same kind of poles as in Group Division}

Your flying carpet has started malfunctioning! As long as it is empty, everything is fine, and it can
fly anywhere in the plane without problems. But as soon as it is loaded with something, such as magical oil lamps,
it can no longer turn on its own, but must turn at so-called spinning poles. Since
the spinning poles can only rotate counterclockwise 90 degrees at a time, the carpet can also only turn exactly
90 degrees counterclockwise at each turn. Moreover, it overheats so quickly that it leaves a burning
trail behind it as soon as it is loaded with something. Therefore, it also cannot turn multiple times at
the same pole, as that would take so long that it would burn up. It also cannot cross its own path or
visit a pole it has already been to.

And of course it is right now, when the carpet is in worse shape than ever, that you need it the most.
Rafaj has scattered all of the sultan's magical oil lamps on the city's spinning poles.
You must now collect as many oil lamps as possible before someone accidentally releases the genies and
disaster strikes.

In the plane there are $n$ spinning poles with integer coordinates and on each pole there is now a magical oil lamp.
You want to visit as many poles as possible with the flying carpet, by moving in the following way:
\begin{itemize}
  \item You may start at any pole, since the carpet can fly freely when empty
  \item You may only move up/down/left/right
  \item You may only turn at poles, and only 90 degrees counterclockwise (so you cannot turn clockwise or reverse direction, but you can continue straight ahead)
  \item Your path may not cross itself, or visit a pole more than once, as the carpet would burn up
\end{itemize}
How many poles can you visit?

\section*{Input}
The first line contains an integer $n$ ($1 \leq n \leq 1000$).
Then follow $n$ lines. The $i$-th line contains two positive integers
$x_{i}, y_{i}$ ($1 \leq x_{i}, y_{i} \leq 10^9$), the coordinates of the $i$-th pole.
No two poles will ever be at the same position.

\section*{Output}
Print a line with an integer, the largest number of poles you can visit.

\section*{Scoring}
Your solution will be tested on a set of test groups.
To earn points for a group, you must pass all test cases in that group.

\noindent
\begin{tabular}{| l | l | l |}
\hline
\textbf{Group} & \textbf{Points} & \textbf{Constraints} \\ \hline
$1$   & $19$       & $n \leq 25, x_{i}, y_{i} \leq 5$ \\ \hline
$2$   & $16$       & $n \leq 80, x_{i}, y_{i} \leq 80$ \\ \hline
$3$   & $20$       & $n \leq 80$ \\ \hline
$4$   & $45$       & No additional constraints. \\ \hline
\end{tabular}
